\begin{proofbox}
\textbf{\textcolor{CProof}{方法一：考场最快法——斜率截距、平方差（优先掌握）}}

\textbf{已知常用公式：}椭圆 $\frac{x^2}{a^2}+\frac{y^2}{b^2}=1$ 的非竖直切线 $y=kx+m$ 满足
\[
\boxed{m^2=a^2k^2+b^2}.
\]
令 $x=-a,a$，两交点的纵坐标为
\[
y_1=m-ak,\quad y_2=m+ak.
\]
于是
\[
\boxed{h_1h_2=|y_1y_2|=|m^2-a^2k^2|=b^2.}
\]
因 $m^2=a^2k^2+b^2>a^2k^2$，故 $|m|>|ak|$，从而 $y_1,y_2$ 与 $m$ 同号。
所以 $h_1+h_2=|y_1|+|y_2|=2|m|$，
\[
S=a(h_1+h_2)=2a|m|\ge2ab.
\]
等号 $\Longleftrightarrow |m|=b\Longleftrightarrow k=0$，即水平切线 $y=\pm b$，切点 $P=(0,\pm b)$。

\textbf{若忘记“切线判定式”，现场两行恢复：}把 $y=kx+m$ 代入椭圆，得到
\[
(b^2+a^2k^2)x^2+2a^2kmx+a^2(m^2-b^2)=0.
\]
相切 $\Longleftrightarrow\Delta=0$，而
\[
\Delta=4a^2b^2\bigl(a^2k^2+b^2-m^2\bigr),
\]
故 $m^2=a^2k^2+b^2$。\textbf{这一步是斜率法的必要依据，不可把任意直线都当成切线。}

\medskip
\textbf{\textcolor{CProof}{方法二：设切点法——完全从熟悉的切线公式出发（保留原主线并纠错）}}

设 $P(x_0,y_0)$ 为椭圆上点，且 $y_0\ne0$，则
\[
\frac{x_0^2}{a^2}+\frac{y_0^2}{b^2}=1,\qquad
\frac{x_0x}{a^2}+\frac{y_0y}{b^2}=1.
\]
令 $x=-a,a$，得
\[
P_1=\left(-a,\frac{b^2}{y_0}\left(1+\frac{x_0}{a}\right)\right),\quad
P_2=\left(a,\frac{b^2}{y_0}\left(1-\frac{x_0}{a}\right)\right).
\]
由于 $-a<x_0<a$，两括号都大于 $0$；但 $y_0$ 可以是负数，\textbf{不能直接删去距离的绝对值}。因此
\[
\begin{aligned}
h_1&=\frac{b^2}{|y_0|}\left(1+\frac{x_0}{a}\right),\\
h_2&=\frac{b^2}{|y_0|}\left(1-\frac{x_0}{a}\right).
\end{aligned}
\]
两式相乘，利用 $1-x_0^2/a^2=y_0^2/b^2$：
\[
\begin{aligned}
h_1h_2
&=\frac{b^4}{y_0^2}\left(1-\frac{x_0^2}{a^2}\right)
=\frac{b^4}{y_0^2}\cdot\frac{y_0^2}{b^2}
=\boxed{b^2}.
\end{aligned}
\]
两式相加：
\[
h_1+h_2=\frac{2b^2}{|y_0|},\qquad
\boxed{S=\frac{2ab^2}{|y_0|}}.
\]
椭圆上 $0<|y_0|\le b$，故 $S\ge2ab$，等号恰在 $y_0=\pm b$、$x_0=0$ 时取得。

\textbf{方法二的速记：} \textbf{令 $x=\pm a$，出现 $1\pm x_0/a$；相乘是平方差，相加是 $2$。}

\medskip
\textbf{\textcolor{CProof}{方法三：三角参数法——半角正切互为倒数}}

写 $P=(a\cos\theta,b\sin\theta)$，$\sin\theta\ne0$，代入切线公式：
\[
\frac{x\cos\theta}{a}+\frac{y\sin\theta}{b}=1.
\]
令 $x=-a,a$ 得
\[
\begin{aligned}
y_1&=b\frac{1+\cos\theta}{\sin\theta}
=b\cot\frac{\theta}{2},\\
y_2&=b\frac{1-\cos\theta}{\sin\theta}
=b\tan\frac{\theta}{2}.
\end{aligned}
\]
因 $\tan(\theta/2)\cot(\theta/2)=1$，有
\[
h_1h_2=|y_1y_2|=b^2.
\]
又 $S=ab\bigl(|\tan(\theta/2)|+|\cot(\theta/2)|\bigr)\ge2ab$；当且仅当两项都为 $1$ 时等号成立，即 $\theta=\frac\pi2$ 或 $\frac{3\pi}{2}$（模 $2\pi$），对应 $P=(0,\pm b)$。

\medskip
\textbf{\textcolor{CProof}{方法四：纯圆几何 + 横向拉伸——理解乘积为什么会出现}}

\textbf{第一步：先证圆。}考虑圆 $u^2+v^2=b^2$，端点 $D_1(-b,0),D_2(b,0)$，其竖直切线 $u=\pm b$。圆上除 $D_1,D_2$ 外的动点 $T$ 的切线交这两条竖线于 $Q_1,Q_2$。先取上半圆情形，记 $s=|D_1Q_1|>0$、$t=|D_2Q_2|>0$。

同一外点引圆的两条切线，切线段相等，所以
\[
|Q_1T|=|Q_1D_1|=s,\qquad
|Q_2T|=|Q_2D_2|=t.
\]
因为 $T$ 位于线段 $Q_1Q_2$ 内部（其横坐标在 $-b$ 与 $b$ 之间），所以 $|Q_1Q_2|=s+t$。另一方面，$Q_1Q_2$ 的水平位移是 $2b$，竖直位移是 $t-s$，由勾股定理
\[
(s+t)^2=|Q_1Q_2|^2=(2b)^2+(t-s)^2.
\]
展开、相消，得到 $4st=4b^2$，即 $st=b^2$。下半圆关于横轴对称，结论相同。

\textbf{第二步：再把圆拉成椭圆。}实施线性变换
\[
(u,v)\mapsto(x,y)=\left(\frac ab u,v\right).
\]
圆变成椭圆，$D_1,D_2,Q_1,Q_2,T$ 依次变成 $A_1,A_2,P_1,P_2,P$，直线和相切关系保持。\textbf{所有竖直线段长度保持不变}，于是
\[
|A_1P_1|\,|A_2P_2|=st=b^2.
\]
圆的对应梯形面积 $\ge2b^2$，横向拉伸后面积乘以 $a/b$，故椭圆情形 $S\ge2ab$。等号对应圆的上下端点，变换后仍为 $P=(0,\pm b)$。

\textbf{说明：}“从外点作两条切线，切线段等长”只在\textbf{圆的证明步骤}中使用；横向拉伸后一般不再等长，因此不可直接用于椭圆。

\medskip
\textbf{\textcolor{CProof}{方法五：只求面积时的最省步骤——定积 + 基本不等式}}

若前面已经求出 $h_1h_2=b^2$，无需再做坐标运算。因为
\[
S=a(h_1+h_2)\ge2a\sqrt{h_1h_2}=2ab.
\]
取等 $h_1=h_2=b$，对应动切线水平，$P=(0,\pm b)$。

\medskip
\textbf{\textcolor{CProof}{短轴端点版本（交换 $x,y$）}}

把端点改成 $B_1(0,-b),B_2(0,b)$，并把交点到端点的\textbf{水平距离}记为 $w_1,w_2$，要求 $x_0\ne0$。由切线公式令 $y=\pm b$，得
\[
w_1=\frac{a^2}{|x_0|}\left(1+\frac{y_0}{b}\right),\qquad
w_2=\frac{a^2}{|x_0|}\left(1-\frac{y_0}{b}\right),
\]
故 $w_1w_2=a^2$，面积 $S=b(w_1+w_2)\ge2ab$，等号 $P=(\pm a,0)$。

\medskip
\textbf{\textcolor{CProof}{本题证明路线选择}}

\textbf{选填题：}优先方法一（已知判定式时几步完成）。\quad
\textbf{不会判定式：}方法二（设点、代 $\pm a$、平方差）。\quad
\textbf{理解和记忆：}方法三（正切与余切互倒）或方法四（圆的切线段 + 拉伸）。
\end{proofbox}
